\documentclass[12pt,a4paper]{article}
\usepackage[margin=25mm, headheight=15pt, headsep=10pt]{geometry}
\usepackage{fontspec}
\usepackage[table]{xcolor} % Note: table option is required for rowcolors
\usepackage{titlesec}
\usepackage{tocloft}
\usepackage{fancyhdr}
\usepackage{booktabs}
\usepackage{tcolorbox}
\tcbuselibrary{skins, breakable}
\usepackage[colorlinks=true, linkcolor=emerald, urlcolor=emerald, bookmarks=true]{hyperref}

\definecolor{emerald}{RGB}{0,102,51}
\definecolor{darkred}{RGB}{139,0,0}
\definecolor{lightgray}{RGB}{245,245,245}

\usepackage[nil]{babel}
\babelprovide[import, main]{bengali}
\babelprovide[import]{arabic}
\babelfont{rm}[Renderer=HarfBuzz,Scale=1.0]{Noto Serif Bengali}
\babelfont{sf}[Renderer=HarfBuzz]{Noto Sans Bengali}
\babelfont{tt}[Renderer=HarfBuzz]{Noto Sans Bengali}
\babelfont[arabic]{rm}[Renderer=HarfBuzz,Scale=1.1]{Noto Naskh Arabic}

% Section styling
\titleformat{\section}{\Large\bfseries\color{emerald}}{\thesection.}{0.5em}{}
\titleformat{\subsection}{\large\bfseries\color{emerald!85!black}}{\thesubsection.}{0.5em}{}
\titleformat{\subsubsection}{\normalsize\bfseries\color{emerald!70!black}}{\thesubsubsection.}{0.5em}{}
\titlespacing*{\section}{0pt}{18pt}{8pt}

% Fancy Header/Footer setup
\pagestyle{fancy}
\fancyhf{} % clear all header and footer fields
\fancyhead[L]{\color{emerald}\small\textbf{\nouppercase{\leftmark}}}
\fancyfoot[L]{\scriptsize\color{black!50}{\lat{AI}}-সংকলিত, আলেম-যাচাইকৃত নয়}
\fancyfoot[C]{\color{emerald!70!black}\thepage}
\renewcommand{\headrulewidth}{0.4pt}
\renewcommand{\headrule}{\hbox to\headwidth{\color{emerald}\leaders\hrule height \headrulewidth\hfill}}
\renewcommand{\footrulewidth}{0pt}

% Custom boxes
% 1. Ayat/Hadith Box
\newtcolorbox{ayatbox}[1][]{
    enhanced, breakable, 
    colback=emerald!5, 
    colframe=emerald, 
    boxrule=0.5pt, 
    arc=3pt, 
    left=10pt, right=10pt, top=10pt, bottom=10pt,
    #1
}

% 2. Quote/Translation Box
\newtcolorbox{quotebox}[1][]{
    enhanced, breakable,
    frame hidden,
    colback=lightgray,
    borderline west={3pt}{0pt}{emerald},
    left=15pt, right=10pt, top=10pt, bottom=10pt,
    sharp corners,
    #1
}

% 3. AI-generated notice box
\newtcolorbox{warnbox}[1][]{
    enhanced, breakable,
    colback=red!4,
    colframe=darkred,
    boxrule=0.8pt,
    arc=3pt,
    left=10pt, right=10pt, top=10pt, bottom=10pt,
    #1
}

% Commands
\newcommand{\arab}[1]{\foreignlanguage{arabic}{#1}}
\newcommand{\kib}[1]{\textcolor{emerald}{\textbf{#1}}}

% Latin (English) fallback: Noto Serif Bengali has NO Latin glyphs
\newfontfamily\latfont[Renderer=HarfBuzz]{Noto Serif}
\newcommand{\lat}[1]{{\latfont #1}}

\setlength{\parskip}{5pt}
\setlength{\parindent}{0pt}

% Fix for missing bullet in Noto Serif Bengali
\renewcommand{\labelitemi}{$\bullet$}



\begin{document}
\begin{center}{\Large\bfseries\color{emerald} \lat{08}-সিজদা-তাসবিহ-দোয়া}\end{center}
\vspace{1em}
\section*{০৮ — সিজদা: তাসবিহ ও দোয়া — \textit{\lat{Subhana} \lat{Rabbiyal} \lat{A'la}}}

\begin{warnbox}
 \textbf{গুরুত্বপূর্ণ নোটিশ:} এই কন্টেন্টটি \lat{AI} দিয়ে প্রস্তুত। কেবল সহীহ/হাসান হাদিস রাখা হয়েছে। আলেম-যাচাইকৃত নয়।
\end{warnbox}

\medskip\hrule\medskip

\subsection*{১. সূত্র কার্ড}

\begin{table}[h]\centering\small
\begin{tabular}{ll}
বিষয় & বিবরণ \\
\hline
\textbf{বিষয়} & সিজদায় যাওয়ার তাকবীর, সিজদায় তাসবিহ, সিজদার দোয়া \\
\textbf{সূত্র} & সহীহ মুসলিম ৪৮২, সুনানে আবু দাউদ ৮৭০ \\
\textbf{মর্যাদা} & \textbf{সুন্নাতে মুআক্কাদাহ} — সিজদায় তাসবিহ ও দোয়া \\
\hline
\end{tabular}
\end{table}

\medskip\hrule\medskip

\subsection*{২. মূল আরবি}

\textbf{সিজদায় তাসবিহ (কমপক্ষে ৩ বার):}
\begin{center}\arab{سُبْحَانَ رَبِّيَ الْأَعْلَىٰ}\end{center}

\textbf{সিজদার দোয়া:}
\begin{center}\arab{سُبْحَانَكَ اللَّهُمَّ رَبَّنَا وَبِحَمْدِكَ اللَّهُمَّ اغْفِرْ لِي}\end{center}

\medskip\hrule\medskip

\subsection*{৩. বাংলা উচ্চারণ}

\textbf{তাসবিহ:} সুবহানা রাব্বিয়াল আ'লা
\textbf{দোয়া:} সুবহানাকাল্লাহুম্মা রাব্বানা ওয়া বিহামদিকা, আল্লাহুম্মাগফির লী

\medskip\hrule\medskip

\subsection*{৪. শব্দ-শব্দ অর্থ}

\begin{table}[h]\centering\small
\begin{tabular}{lll}
আরবি & বাংলা & \lat{English} \\
\hline
\textbf{\arab{سُبْحَانَ}} & পবিত্র ঘোষণা & \textit{\lat{Glorified} \lat{is}} \\
\textbf{\arab{رَبِّيَ}} & আমার রব & \textit{\lat{my} \lat{Lord}} \\
\textbf{\arab{الْأَعْلَىٰ}} & সর্বোচ্চ & \textit{\lat{the} \lat{Most} \lat{High}} \\
\textbf{\arab{سُبْحَانَكَ}} & তোমাকে পবিত্র ঘোষণা & \textit{\lat{Glorified} \lat{be} \lat{You}} \\
\textbf{\arab{رَبَّنَا}} & আমাদের রব & \textit{\lat{Our} \lat{Lord}} \\
\textbf{\arab{وَبِحَمْدِكَ}} & ও তোমার প্রশংসায় & \textit{\lat{and} \lat{with} \lat{Your} \lat{praise}} \\
\textbf{\arab{اغْفِرْ} \arab{لِي}} & আমাকে ক্ষমা কর & \textit{\lat{forgive} \lat{me}} \\
\hline
\end{tabular}
\end{table}

\medskip\hrule\medskip

\subsection*{৫. পূর্ণ বাংলা অনুবাদ}

\textbf{\lat{“}আমার রব সর্বোচ্চের কাছে পবিত্রতা।\lat{”}}

\textbf{\lat{“}হে আল্লাহ! আমি তোমাকে পবিত্র ঘোষণা করি — আমাদের রব তুমি, তোমার প্রশংসায়। হে আল্লাহ! আমাকে ক্ষমা কর।\lat{”}}

\medskip\hrule\medskip

\subsection*{৬. সংক্ষিপ্ত ব্যাখ্যা}

\subsubsection*{\textbf{সিজদা কী?}}
\lat{Sujud} মানে মাথা, নাক, হাত, হাঁটু, পায়ের আঙুল মাটিতে রেখে আল্লাহর কাছে আত্মসমর্পণ। হাদিসে এসেছে: \textbf{\lat{“}বান্দা তার রবের সবচেয়ে কাছে থাকে সিজদার অবস্থায়।\lat{”}} (সহীহ মুসলিম ৪৮২)।

\subsubsection*{\textbf{"আ'লা" — সর্বোচ্চ}}
\lat{Al-A'la} মানে সর্বোচ্চ। রুকুতে "আযীম" (মহান), সিজদায় "আ'লা" (সর্বোচ্চ) — সিজদায় বান্দা সবচেয়ে নিচে, কিন্তু আল্লাহ সবচেয়ে উপরে।

\subsubsection*{\textbf{৭টি অঙ্গ মাটিতে}}
নাক, দুই হাতের হাতল, দুই হাঁটু, পায়ের আঙুল — এগুলো অবশ্যই মাটিতে স্পর্শ করতে হবে।

\medskip\hrule\medskip

\subsection*{৭. খুশুর জন্য মূল পয়েন্টস}

\begin{table}[h]\centering\small
\begin{tabular}{ll}
পয়েন্ট & কীভাবে প্রয়োগ করবেন \\
\hline
\textbf{সময়কাল} & প্রতিটি সিজদায় কমপক্ষে ৩ বার তাসবিহ বলুন \\
\textbf{দোয়া} & সিজদায় আল্লাহর কাছে যা চান তা চান — সবচেয়ে কাছের মুহূর্ত \\
\textbf{বিরতি} & দুই সিজদার মাঝে "রাব্বিগফির লী ওয়ারহামনী" বলুন \\
\textbf{নজর} & চোখ সিজদার জায়গায় — কিবলার দিকে \\
\hline
\end{tabular}
\end{table}

\medskip\hrule\medskip

\subsection*{৮. সংশ্লিষ্ট হাদিস}

\begin{table}[h]\centering\small
\begin{tabular}{ll}
হাদিস & গ্রন্থ \\
\hline
\textbf{\lat{“}বান্দা তার রবের সবচেয়ে কাছে থাকে সিজদার অবস্থায় — তাই সিজদায় বেশি বেশি দোয়া করো।\lat{”}} & সহীহ মুসলিম ৪৮২ \\
\textbf{\lat{“}সিজদায় 'সুবহানা রাব্বিয়াল আ'লা' ৩ বার বলো।\lat{”}} & সুনানে আবু দাউদ ৮৭০ \\
\hline
\end{tabular}
\end{table}

\medskip\hrule\medskip

\textit{শেষ আপডেট: আগস্ট ২০২৬}

\end{document}
